% !TEX root = ../MLEvolve.tex

\section{Related Work}
\label{related_works}

\subsection{Automated Machine Learning Algorithm Discovery}
To address the unique challenges of MLE, a dedicated class of coding agents has been developed~\citep{aide,mle-star,mlzero,ml-master}, with many evaluated on benchmarks such as MLE-Bench~\citep{mle-bench}. These agents primarily frame the problem as a search for an optimal code-based solution. Early works like AIDE~\citep{aide} employ greedy search, which is susceptible to local optima. Subsequent frameworks adopt more structured exploration. ML-Master~\citep{ml-master} and AIRA-Dojo~\citep{dojo} use MCTS, MARS~\citep{chen2026mars} introduces budget-aware MCTS with contrastive reflection, and FM-Agent~\citep{li2025fm} applies evolutionary multi-island parallel search. Other works explore agent collaboration, such as R\&D-Agent~\citep{rdagent} with researcher-developer combination and AIBuildAI~\citep{zhang2026aibuildai} with hierarchical multi-agent coordination. Several methods also incorporate external knowledge or memory. AutoMind~\citep{automind} and Leeroo~\citep{nadafian2026kapso} ground search with domain knowledge bases, while ML-Master 2.0~\citep{zhu2026mlmaster2} introduces hierarchical cognitive caching for cross-task knowledge distillation. However, these methods commonly suffer from inter-branch information isolation and the inability to accumulate and reuse experience from past trials. Our method addresses these limitations from a self-evolving perspective, enabling the agent to continuously adapt its search behavior, accumulate experience, and refine solutions during long-horizon optimization.


\subsection{Graph-based Planning and Search}
Early methods that combine graph structures with MCTS, often referred to as MCGS~\citep{mcgs1,mcgs2}, were primarily developed for planning and reinforcement learning tasks with well-defined state spaces, where identical states are merged to compress the search space. Recent graph-based frameworks such as LocAgent~\citep{locagent2025} and CodexGraph~\citep{codexgraph2025} use graphs as static dependency representations for retrieval or localization, but these graphs do not evolve during search. In contrast, our MCGS targets open-ended LLM-based code generation, where each node represents a distinct candidate solution. The graph structure is not used to compress the state space, but to enable cross-branch information flow, trajectory reuse, and solution composition through dynamic reference edges.

\subsection{Memory and Experience Mechanisms for LLM Agents}
Memory mechanisms have been explored to improve LLM agent performance in iterative tasks~\citep{zhang2025memsurvey}. Recent work on long-term episodic memory~\citep{xu2026amem} enables agents to accumulate and retrieve experiential records across extended horizons, supporting more informed subsequent decisions. In the MLE domain, recent works further explore experience reuse. ROME~\citep{zhang2026rome} introduces ``reasoning gradients'' as structured optimization directions and stores successful trajectories as momentum memory, MARS~\citep{chen2026mars} extracts insights through contrastive reflection over historical attempts, and ML-Master 2.0~\citep{zhu2026mlmaster2} introduces hierarchical cognitive caching for cross-task knowledge distillation. While these methods advance experience reuse, most require additional LLMs for reflection or summarization. Our retrospective memory automatically accumulates and retrieves experience 
without requiring additional LLMs for explicit reflection, and further incorporates a static domain knowledge base for cold-start initialization.

